% Source: Weinberg GR, physical PDF pp. 116--118 (printed pp. 91--93).
% Coverage: Section 4.1, The Principle of General Covariance; no numbered
%           equations.
% Figures/tables/footnotes: four linked reference markers; no figures or tables.
% Status: source-reviewed and compile-clean.
% Uncertainties: none.

\subsection{The Principle of General Covariance}\label{sec:4.1}

In the last chapter we demonstrated one way of using the Principle of
Equivalence to assess the effects of gravitation on physical systems: We wrote
down the equations that hold, for general gravitational fields, in locally
inertial coordinate systems (i.e., the equations of special relativity, such as
\(\dd^2\xi^\alpha/\dd\tau^2=0\)) and then performed a coordinate
transformation to find the corresponding equations in the laboratory
coordinate system. We could continue to follow this approach, but it would
lead us into very tedious calculations when we come to the field equations for
electromagnetism and gravitation. Instead, we shall follow a different method,
one that is of precisely the same physical content, but is much more elegant in
appearance and convenient in execution. This method is based on an alternative
version of the Principle of Equivalence, known as the Principle of General
Covariance. It states that a physical equation holds in a general
gravitational field, if two conditions are met:

\begin{enumerate}
  \item The equation holds in the absence of gravitation; that is, it agrees
  with the laws of special relativity when
  \(g_{\mu\nu}=\eta_{\mu\nu}=\diag(-1,+1,+1,+1)\) and when the affine
  connection \(\Gamma^\rho{}_{\mu\nu}\) vanishes.

  \item The equation is generally covariant; that is, it preserves its form
  under a general coordinate transformation \(x^\mu\to x'^\mu\).
\end{enumerate}

To see that the Principle of General Covariance follows from the Principle of
Equivalence, let us suppose that we are in an arbitrary gravitational field,
and consider any equation that satisfies the two above conditions. From (2),
we learn that the equation will be true in all coordinate systems if it is
true in any one coordinate system. But at any given point there is a class of
coordinate systems, the locally inertial systems, in which the effects of
gravitation are absent. Condition (1) then tells us that our equation holds in
these systems, and hence in all other coordinate systems.

It should be stressed that general covariance by itself is empty of physical
content.~\hyperref[ch4-ref-1]{[1]} Any equation can be made generally
covariant by writing it in any one coordinate system, and then working out
what it looks like in other arbitrary coordinate systems. Indeed, from
childhood we have become familiar with the appearance of physical equations
in non-Cartesian systems, such as polar coordinates, and in noninertial
systems such as rotating coordinates. The significance of the Principle of
General Covariance lies in its statement about the effects of gravitation:
that a physical equation by virtue of its general covariance will be true in
a gravitational field if it is true in the absence of gravitation.

The meaning of general covariance can be brought forward by comparing it with
Lorentz invariance. Just as any equation can be made generally covariant, so
any equation can be made Lorentz invariant, by writing it in one coordinate
system and then working out what it looks like after a Lorentz transformation.
However, if we do this with a nonrelativistic equation like Newton's second
law, we find after making it Lorentz invariant that a new quantity has entered
the equation, which of course is the velocity of the coordinate frame with
respect to the original reference frame. The requirement that this velocity
not appear in the transformed equation is what we call the Principle of
Special Relativity, or ``Lorentz invariance'' for short, and this requirement
places very powerful restrictions on the original equation. Similarly, when
we make an equation generally covariant, new ingredients will enter, that is,
the metric tensor \(g_{\mu\nu}\) and the affine connection
\(\Gamma^\rho{}_{\mu\nu}\). The difference is that we do not require that
these quantities drop out at the end, and hence we do not obtain any
restrictions on the equation we start with; rather, we exploit the presence of
\(g_{\mu\nu}\) and \(\Gamma^\rho{}_{\mu\nu}\) to represent gravitational
fields.

To put this briefly: The Principle of General Covariance is not an invariance
principle, like the Principle of Galilean or Special Relativity, but is instead
a statement about the effects of gravitation, and about nothing else. In
particular, general covariance does not imply Lorentz invariance---there are
generally covariant theories of gravitation that allow the construction of
inertial frames at any point in a gravitational field, but that satisfy
Galilean relativity rather than special relativity in these
frames.~\hyperref[ch4-ref-1a]{[1a]}

Any physical principle, such as general covariance, which takes the form of an
invariance principle but whose content is actually limited to a restriction
on the interactions of one particular field, is called a \textit{dynamic
symmetry}.~\hyperref[ch4-ref-2]{[2]} There are other dynamic symmetries of
importance in physics, such as local gauge invariance, which governs the
interactions of the electromagnetic field, and chiral symmetry, which governs
the interactions of the pi-meson field.~\hyperref[ch4-ref-3]{[3]} We shall
return to the analogy between general relativity and electrodynamics several
times in following chapters.

The Principle of General Covariance can only be applied on a scale that is
small compared with the spacetime distances typical of the gravitational
field, for it is only on this small scale that we are assured by the Principle
of Equivalence of being able to construct a coordinate system in which the
effects of gravitation are absent. For instance, the radius of the moon is not
so very much smaller than the earth--moon separation, so we cannot accurately
calculate the motion of the moon by finding generally covariant equations that
reduce to the correct equations for a freely moving moon in the absence of
gravitation. We can, however, treat the moon as a ball of rock and calculate
its motion by applying the Principle of General Covariance to determine the
gravitational force on each infinitesimal element of the lunar mass.

There are in general many generally covariant equations that reduce to a given
special-relativistic equation in the absence of gravitation. However, because
we only apply the Principle of General Covariance on a small scale compared
with the scale of the gravitational field, we usually expect that it is only
\(g_{\mu\nu}\) and its first derivatives that enter our generally covariant
equations. With this understanding we shall see in this and the next chapter
that the Principle of General Covariance makes an unambiguous statement about
the effects of gravitational fields on any system, or part of a system, that
is sufficiently small.
